The course contributes to the following three exit qualifications of the Research Master's programme in Communication Science:

\textit{Expertise in empirical research}

\begin{enumerate}
    \setcounter{enumi}{2}
    \item Knowledge and Understanding: Have in-depth knowledge and a thorough understanding of advanced research designs and methods.
    \item Skills and abilities: Are able, independently and on their own, to set up, conduct, report, and interpret advanced academic research.
\end{enumerate}

\textit{Academic abilities and attitudes}

\begin{enumerate}
    \setcounter{enumi}{5}
    \item Academic attitudes: Accept that scientific knowledge is always `work in progress' and that arguments must be considered and conclusions drawn on the basis of empirical results and valid criticism.
\end{enumerate}

The course contributes to these exit qualifications at the level of \emph{fundamental} techniques of Big Data and Automated Content Analysis. In particular, students learn to justify methodological choices, acknowledge limitations and uncertainties in their own research, and respond constructively to critical questions. These academic attitudes are assessed through the individual project pitch and the final individual project.

The exit qualifications are elaborated in the following specifications:

\begin{enumerate}
    \setcounter{enumi}{2}
    \item Knowledge and Understanding: Have in-depth knowledge and a thorough understanding of advanced research designs and methods.
    \begin{enumerate}
        \item[3.1.] Have in-depth knowledge and a thorough understanding of advanced research designs and methods, including their value and limitations.
        \item[3.2.] Have in-depth knowledge and a thorough understanding of advanced techniques for data analysis.
    \end{enumerate}

    \item Skills and abilities: Are able, independently and on their own, to set up, conduct, report, and interpret advanced academic research.
    \begin{enumerate}
        \item[4.1.] Are able to formulate research questions and hypotheses for advanced empirical studies.
        \item[4.2.] Are able to develop a research plan, choose appropriate and suitable research designs and methods for advanced empirical studies, and justify the underlying choices.
        \item[4.3.] Are able to assess the validity and reliability of advanced empirical research, and to judge the scientific and professional value of findings from advanced empirical research.
        \item[4.4.] Are able to apply advanced empirical research methods.
    \end{enumerate}

    \item Academic attitudes.
    \begin{enumerate}
        \item[6.1.] Regularly assess their own assumptions, strengths, and weaknesses critically.
        \item[6.2.] Accept that scientific knowledge is always `work in progress' and that something regarded as `true' may be proven to be false, and vice versa.
        \item[6.3.] Are keen to acquire new knowledge, skills, and abilities.
        \item[6.4.] Are willing to share and discuss arguments, results, and conclusions, including submitting one's own work to peer review.
        \item[6.5.] Are convinced that academic debates should not be conducted on the basis of rhetorical qualities but that arguments must be considered and conclusions drawn on the basis of empirical results and valid criticism.
    \end{enumerate}
\end{enumerate}
